% Zander Carpio, curriculum vitae.
% Build:  tectonic cv/cv.tex --outdir assets/cv
\documentclass[10pt,a4paper]{article}

\usepackage[T1]{fontenc}
\usepackage[utf8]{inputenc}
\usepackage[margin=0.6in]{geometry}
\usepackage{titlesec}
\usepackage{enumitem}
\usepackage{xcolor}
\usepackage{array}
\usepackage[hidelinks]{hyperref}

\definecolor{ink}{HTML}{111C2E}
\definecolor{accent}{HTML}{C2410C}
\definecolor{muted}{HTML}{56657A}

\color{ink}
\pagestyle{empty}
\setlength{\parindent}{0pt}
\linespread{1.0}

% Section headings: small caps, accent rule underneath
\titleformat{\section}
  {\normalfont\normalsize\bfseries\scshape\color{accent}}
  {}{0pt}{}[\vspace{-7pt}\textcolor{accent!35}{\rule{\linewidth}{0.6pt}}]
\titlespacing{\section}{0pt}{9pt}{4pt}

% One entry: title, org, dates
\newcommand{\entry}[3]{%
  \textbf{#1}\hfill{\small\textcolor{muted}{#3}}\\[1pt]
  {\small\textcolor{accent}{#2}}\\[2pt]
}

\newcommand{\proj}[3]{%
  \textbf{#1}\ \ {\small\textcolor{muted}{#2}}\hfill{\small\textcolor{muted}{#3}}\\[1pt]
}

\begin{document}

% ------------------------------------------------------------------ header
{\LARGE\bfseries Zander Carpio}\\[5pt]
{\small\textcolor{muted}{
  Baguio City, Cordillera Administrative Region, Philippines \ \textbar\
  \href{mailto:zandercarpio.ph@gmail.com}{zandercarpio.ph@gmail.com}\\
  \href{https://www.linkedin.com/in/zander-carpio/}{linkedin.com/in/zander-carpio}
  \ \textbar\
  \href{https://github.com/zandercarpioph-create}{github.com/zandercarpioph-create}
}}

\vspace{9pt}
{\small
Mathematics undergraduate working on quantitative and analytical problems:
epidemic dynamics, geospatial disease surveillance, official statistics, and
computational number theory. Seeking analytical, data and research roles.
}

% --------------------------------------------------------------- education
\section{Education}

\entry{BS Mathematics}{University of the Philippines Baguio}{Aug 2023 to 2027 (expected)}
{\small
\textbf{DOST-SEI Undergraduate Scholar}, a competitive national merit scholarship
from the Philippine Department of Science and Technology. Undergraduate thesis in
number theory.\\[2pt]
\textit{Relevant coursework:} Real Analysis, Complex Analysis, Analytic Number
Theory, Operations Research, Physics, Accounting.
}

\vspace{6pt}
\entry{Secondary Education}{Aurora National Science High School}{2017 to 2023}

% -------------------------------------------------------------- experience
\section{Experience}

\entry{Intern, Statistical Operations and Coordination Division}{Philippine Statistics Authority, Baguio City}{Jun to Jul 2026}
{\small
\begin{itemize}[leftmargin=12pt,itemsep=1pt,topsep=1pt]
  \item Statistical operations, data visualisation and coordination work inside
        the national statistics agency responsible for official statistics, the
        census and civil registration.
  \item Worked with survey and administrative data in its delivered state,
        including inconsistent labelling and failed joins.
\end{itemize}
}

\vspace{3pt}
\entry{Peer Mentor, Learning Resource Center}{University of the Philippines Baguio, Contract}{Sep 2025 to present}
{\small Mentoring undergraduates through mathematics coursework.}

\vspace{3pt}
\entry{Mathematics Tutor}{Freelance, Hybrid}{Jul 2021 to Jul 2024}
{\small Three years tutoring mathematics at secondary and lower undergraduate level.}

\vspace{3pt}
\entry{Social Media Manager}{Freelance, Remote}{Jul 2022 to present}
{\small Four years running social media for small businesses alongside a full-time degree.
Campaign management, scheduling and reporting in GoHighLevel and SocialPilot, with graphic
design and content produced in house. Client names withheld.}

% ---------------------------------------------------------------- projects
\section{Quantitative Projects}

{\small
\proj{Dengue surveillance analysis, Baguio City}{R, QGIS}{Jan 2024}
Reconstructed barangay-level surveillance across 127 barangays and 23 reporting
weeks, joined on PSGC codes and mapped as a choropleth. Showed the burden is
concentrated rather than citywide, the worst tenth of barangays carrying 32\% of
cases, and confirmed this was not an artefact of barangay size by re-deriving it
from population-normalised incidence (32.4\% against 31.7\%). That check exposed a
denominator failure in the commercial core, where residential population
understates the at-risk daytime population enough to yield nominal rates above
4,000 per 1,000.

\vspace{4pt}
\proj{Compartmental epidemic models}{R, deSolve, Shiny}{May 2026}
SIR and SEIR systems with mass-action incidence, integrated with \texttt{deSolve}
and published as interactive Shiny applications in which the parameters can be
varied and the epidemic curve re-solved.

\vspace{4pt}
\proj{Cholera transmission with an environmental reservoir}{R, MATLAB}{May 2026}
SEIRB model coupling direct transmission to a bacterial reservoir, built as a
797-line Shiny application with fourteen tunable parameters and cross-implemented
in MATLAB.

\vspace{4pt}
\proj{Kaprekar triples, undergraduate thesis}{PARI/GP, LaTeX}{2026, ongoing}
Extending Iannucci and Foster (2005). Reduced an infinite search to a finite one by
proving every triple is a root of a congruence, bounding the count at four per power
of ten. Worked to a rule that no result counts until two independently written
programs agree, which exposed a counting error in earlier code.
}

% ------------------------------------------------------------------ skills
\section{Technical Skills}

{\small
\begin{tabular}{@{}>{\bfseries}p{2.15cm}p{15.3cm}@{}}
Languages & R, Python, MATLAB, C, PARI/GP \\[2pt]
Libraries & deSolve, Shiny, ggplot2, dplyr, tidyr, NumPy \\[2pt]
Tools     & QGIS, LaTeX, Git, Excel, GoHighLevel, SocialPilot \\
\end{tabular}
}

\end{document}
